% Logros del laboratorio 5 frente a sus objetivos y trabajo del docente.
% Generado por ../generar_logros.py desde ../logros_datos.py: editar los datos y regenerar.
\documentclass[a4paper,11pt]{article}
% Preámbulo común de los informes de laboratorio (II/2026). Se incluye desde laboratorios/labNN_*/
%   % Preámbulo común de los informes de laboratorio (II/2026). Se incluye desde laboratorios/labNN_*/
%   % Preámbulo común de los informes de laboratorio (II/2026). Se incluye desde laboratorios/labNN_*/
%   \input{../informe_comun}
\input{../../gemelo-digital/manual/estilo_manual}
\input{../../gemelo-digital/manual/tecnico/macros}
\graphicspath{{../../documentacion/planos_isa/pid/}{../../documentacion/planos_isa/electrico/}{../../documentacion/manuales_inversa/}
  {../../database/mockups/}{../../database/documentacion/}{../../database/modelo/}{../../database/interfaz/}
  {../../gemelo-digital/manual/img/}{../../modelo-3d/informe/figuras/}{../../modelo-3d/levantamiento/}{../../../../global/fotos-planta/}{../figuras/}}
\usepackage{pdflscape}
\lstset{showstringspaces=false}
% plano A3 en una página apaisada: \plano[página]{archivo}{leyenda}
\newcommand{\plano}[3][1]{\begin{landscape}\begin{figure}[p]\centering\includegraphics[page=#1, width=\linewidth, height=0.86\textwidth, keepaspectratio]{#2}\caption{#3}\end{figure}\end{landscape}}
\setlength{\headheight}{14pt}
\newcommand{\nlab}{}
\newcommand{\tlab}{}
\lhead{\small Brazo XK-JS1-1 · Informe del laboratorio \nlab}
\rhead{\small UMSA · II/2026}

% \portada{número}{título}{inicio}{entrega}{semanas}
\newcommand{\portada}[5]{%
\renewcommand{\nlab}{#1}%
\begin{titlepage}
\centering
\vspace*{0.6cm}
{\large Universidad Mayor de San Andrés · Facultad de Ingeniería\par}
{\large Carrera de Ingeniería Electrónica\par}
\vspace{0.8cm}
{\color{naranja!80!black}\rule{\linewidth}{1.5pt}\par}
\vspace{0.3cm}
{\Large Informe del Laboratorio Integrado Nº #1\par}
\vspace{0.2cm}
{\huge\bfseries #2\par}
\vspace{0.3cm}
{\color{naranja!80!black}\rule{\linewidth}{1.5pt}\par}
\vspace{0.8cm}
{\large Planta: brazo robótico XK-JS1-1 de 4 grados de libertad\par}
\vspace{0.3cm}
{\large ETN-902 Sistemas de Control II · ETN-1034 Aplicación de Técnicas de Control · ETN-1000 Bases de Datos\par}
\vspace{1cm}
\begin{tabular}{rl}
\textbf{Docente:} & Ing. Jorge Antonio Nava Amador\\
\textbf{Auxiliares:} & Limber Ajnota · Fabián Plata · Franz Choque\\
\textbf{Inicio de trabajo:} & #3\\
\textbf{Entrega:} & #4 (#5)\\
\end{tabular}
\vspace{1cm}

\begin{tabular}{ll}
\multicolumn{2}{l}{\textbf{Escuadra del brazo de 4 GDL}}\\[2pt]
Control (ETN-902 / ETN-1034): & Gabriel Andres Gomez Calla\\
 & Jose Lenin Mamani Choque\\
 & Miguel Alejandro Lima Huarachi\\
 & Román Abner Machicado Arequipa\\[2pt]
Cómputo (ETN-1000): & Miguel Alejandro Lima Huarachi\\
 & Marco Eduardo Chura Rondo\\
\end{tabular}
\vfill
{\small Repositorio: \url{https://github.com/GabAndGom/Brazo-de-4-grados-de-libertad}\par}
{\small La Paz, Bolivia · gestión II/2026\par}
\end{titlepage}
}

% Tabla de entregables (Definition of Done): \begin{dod} ... \end{dod}, filas: entregable & estado & evidencia\\
\newenvironment{dod}{\begin{center}\small\begin{longtable}{>{\raggedright}p{5.2cm} >{\raggedright}p{1.8cm} >{\raggedright\arraybackslash}p{7.6cm}}
\toprule \textbf{Entregable} & \textbf{Estado} & \textbf{Evidencia / observación}\\ \midrule \endhead}
{\bottomrule\end{longtable}\end{center}}

\input{../../gemelo-digital/manual/estilo_manual}
% Macros del manual técnico: ladder, estados, código y rutas de imágenes
\usetikzlibrary{shapes.multipart}
\graphicspath{{../img/}{../../../../../global/fotos-planta/}{../../../modelo-3d/informe/figuras/}{../../../modelo-3d/levantamiento/}{../../../database/modelo/}{../../../documentacion/manuales_inversa/}{../../../documentacion/planos_isa/electrico/}{../../../database/interfaz/}}

\lstdefinestyle{py}{language=Python, keywordstyle=\color{azulg}\bfseries, commentstyle=\color{grisos}\itshape,
  stringstyle=\color{verde}, basicstyle=\ttfamily\scriptsize, numbers=left, numberstyle=\tiny\color{grisos}}
\lstdefinestyle{js}{language=Java, morekeywords={const,let,function,async,await,return,for,of,if,else}, keywordstyle=\color{azulg}\bfseries,
  commentstyle=\color{grisos}\itshape, stringstyle=\color{verde}, basicstyle=\ttfamily\scriptsize, numbers=left, numberstyle=\tiny\color{grisos}}
\lstdefinestyle{awl}{basicstyle=\ttfamily\tiny, commentstyle=\color{grisos}\itshape, morecomment=[l]{//},
  keywordstyle=\color{azulg}\bfseries, morekeywords={LD,LDN,A,AN,O,ON,EU,ED,S,R,MOVB,MOVW,MOVD,PLS,TON,CALL,CRET,LPS,LRD,LPP,ALD,OLD,FILL,LDW,LDB,LDD,AW,AD,AB,OW,OD},
  numbers=left, numberstyle=\tiny\color{grisos}, frame=single, backgroundcolor=\color{gris}}
\lstdefinestyle{sql}{language=SQL, keywordstyle=\color{azulg}\bfseries, commentstyle=\color{grisos}\itshape,
  stringstyle=\color{verde}, basicstyle=\ttfamily\tiny, numbers=left, numberstyle=\tiny\color{grisos},
  morekeywords={REFERENCES,SERIAL,SMALLSERIAL,BIGSERIAL,TIMESTAMPTZ,BOOLEAN,INET,RETURNS,TRIGGER,LANGUAGE,DECLARE,RAISE,EXCEPTION,VIEW,INDEX}}

\tikzset{
  ld/.style={line width=0.7pt},
  pics/na/.style={code={\fill[white](-0.12,-0.26) rectangle (0.12,0.26);
     \draw[thick](-0.1,-0.24)--(-0.1,0.24)(0.1,-0.24)--(0.1,0.24);
     \node[font=\scriptsize\ttfamily, above] at (0,0.22){#1};}},
  pics/nc/.style={code={\fill[white](-0.12,-0.26) rectangle (0.12,0.26);
     \draw[thick](-0.1,-0.24)--(-0.1,0.24)(0.1,-0.24)--(0.1,0.24);\draw(-0.17,-0.24)--(0.17,0.24);
     \node[font=\scriptsize\ttfamily, above] at (0,0.22){#1};}},
  pics/pf/.style={code={\fill[white](-0.12,-0.26) rectangle (0.12,0.26);
     \draw[thick](-0.1,-0.24)--(-0.1,0.24)(0.1,-0.24)--(0.1,0.24);
     \node[font=\tiny\bfseries] at (0,0){P};}},
  pics/cmp/.style n args={3}{code={\fill[white](-0.12,-0.26) rectangle (0.12,0.26);
     \draw[thick](-0.1,-0.24)--(-0.1,0.24)(0.1,-0.24)--(0.1,0.24);
     \node[font=\tiny\ttfamily] at (0,0){#2};
     \node[font=\scriptsize\ttfamily, above] at (0,0.22){#1};
     \node[font=\scriptsize\ttfamily, below] at (0,-0.22){#3};}},
  pics/bob/.style n args={3}{code={\fill[white](-0.2,-0.26) rectangle (0.2,0.26);
     \draw[thick](-0.05,0.24) arc (110:250:0.255);\draw[thick](0.05,0.24) arc (70:-70:0.255);
     \node[font=\tiny\bfseries] at (0,0){#2};
     \node[font=\scriptsize\ttfamily, above] at (0,0.22){#1};
     \node[font=\scriptsize\ttfamily, below] at (0,-0.22){#3};}},
  fb/.style={draw, fill=white, font=\scriptsize\ttfamily, align=left, inner sep=3pt, anchor=west},
  barra/.style={line width=2pt},
  estado/.style={draw, rounded corners=3pt, fill=white, align=center, font=\scriptsize, minimum width=22mm, minimum height=10mm, inner sep=3pt},
  estadon/.style={estado, fill=naranja!14, draw=naranja!70!black},
  trans/.style={font=\scriptsize, align=center, fill=white, inner sep=1pt},
  lugar/.style={circle, draw, thick, minimum size=7mm, fill=white},
  transicion/.style={rectangle, fill=black, minimum width=1.6mm, minimum height=7mm, inner sep=0},
  tabla/.style={rectangle split, rectangle split parts=2, draw, fill=white, font=\tiny\ttfamily, align=left, inner sep=2pt, rounded corners=1pt, text width=33mm},
}
\newcommand{\pendiente}[1]{\par\smallskip\noindent\colorbox{naranja!12}{\parbox{\dimexpr\linewidth-2\fboxsep}{\textbf{\color{naranja!80!black}Pendiente.} #1}}\par\smallskip}
\newcommand{\cumple}{\textcolor{verde}{\textbf{cumplido}}}
\newcommand{\parcial}{\textcolor{naranja}{\textbf{parcial}}}
\newcommand{\falta}{\textcolor{rojo}{\textbf{pendiente}}}

\graphicspath{{../../documentacion/planos_isa/pid/}{../../documentacion/planos_isa/electrico/}{../../documentacion/manuales_inversa/}
  {../../database/mockups/}{../../database/documentacion/}{../../database/modelo/}{../../database/interfaz/}
  {../../gemelo-digital/manual/img/}{../../modelo-3d/informe/figuras/}{../../modelo-3d/levantamiento/}{../../../../global/fotos-planta/}{../figuras/}}
\usepackage{pdflscape}
\lstset{showstringspaces=false}
% plano A3 en una página apaisada: \plano[página]{archivo}{leyenda}
\newcommand{\plano}[3][1]{\begin{landscape}\begin{figure}[p]\centering\includegraphics[page=#1, width=\linewidth, height=0.86\textwidth, keepaspectratio]{#2}\caption{#3}\end{figure}\end{landscape}}
\setlength{\headheight}{14pt}
\newcommand{\nlab}{}
\newcommand{\tlab}{}
\lhead{\small Brazo XK-JS1-1 · Informe del laboratorio \nlab}
\rhead{\small UMSA · II/2026}

% \portada{número}{título}{inicio}{entrega}{semanas}
\newcommand{\portada}[5]{%
\renewcommand{\nlab}{#1}%
\begin{titlepage}
\centering
\vspace*{0.6cm}
{\large Universidad Mayor de San Andrés · Facultad de Ingeniería\par}
{\large Carrera de Ingeniería Electrónica\par}
\vspace{0.8cm}
{\color{naranja!80!black}\rule{\linewidth}{1.5pt}\par}
\vspace{0.3cm}
{\Large Informe del Laboratorio Integrado Nº #1\par}
\vspace{0.2cm}
{\huge\bfseries #2\par}
\vspace{0.3cm}
{\color{naranja!80!black}\rule{\linewidth}{1.5pt}\par}
\vspace{0.8cm}
{\large Planta: brazo robótico XK-JS1-1 de 4 grados de libertad\par}
\vspace{0.3cm}
{\large ETN-902 Sistemas de Control II · ETN-1034 Aplicación de Técnicas de Control · ETN-1000 Bases de Datos\par}
\vspace{1cm}
\begin{tabular}{rl}
\textbf{Docente:} & Ing. Jorge Antonio Nava Amador\\
\textbf{Auxiliares:} & Limber Ajnota · Fabián Plata · Franz Choque\\
\textbf{Inicio de trabajo:} & #3\\
\textbf{Entrega:} & #4 (#5)\\
\end{tabular}
\vspace{1cm}

\begin{tabular}{ll}
\multicolumn{2}{l}{\textbf{Escuadra del brazo de 4 GDL}}\\[2pt]
Control (ETN-902 / ETN-1034): & Gabriel Andres Gomez Calla\\
 & Jose Lenin Mamani Choque\\
 & Miguel Alejandro Lima Huarachi\\
 & Román Abner Machicado Arequipa\\[2pt]
Cómputo (ETN-1000): & Miguel Alejandro Lima Huarachi\\
 & Marco Eduardo Chura Rondo\\
\end{tabular}
\vfill
{\small Repositorio: \url{https://github.com/GabAndGom/Brazo-de-4-grados-de-libertad}\par}
{\small La Paz, Bolivia · gestión II/2026\par}
\end{titlepage}
}

% Tabla de entregables (Definition of Done): \begin{dod} ... \end{dod}, filas: entregable & estado & evidencia\\
\newenvironment{dod}{\begin{center}\small\begin{longtable}{>{\raggedright}p{5.2cm} >{\raggedright}p{1.8cm} >{\raggedright\arraybackslash}p{7.6cm}}
\toprule \textbf{Entregable} & \textbf{Estado} & \textbf{Evidencia / observación}\\ \midrule \endhead}
{\bottomrule\end{longtable}\end{center}}

\input{../../gemelo-digital/manual/estilo_manual}
% Macros del manual técnico: ladder, estados, código y rutas de imágenes
\usetikzlibrary{shapes.multipart}
\graphicspath{{../img/}{../../../../../global/fotos-planta/}{../../../modelo-3d/informe/figuras/}{../../../modelo-3d/levantamiento/}{../../../database/modelo/}{../../../documentacion/manuales_inversa/}{../../../documentacion/planos_isa/electrico/}{../../../database/interfaz/}}

\lstdefinestyle{py}{language=Python, keywordstyle=\color{azulg}\bfseries, commentstyle=\color{grisos}\itshape,
  stringstyle=\color{verde}, basicstyle=\ttfamily\scriptsize, numbers=left, numberstyle=\tiny\color{grisos}}
\lstdefinestyle{js}{language=Java, morekeywords={const,let,function,async,await,return,for,of,if,else}, keywordstyle=\color{azulg}\bfseries,
  commentstyle=\color{grisos}\itshape, stringstyle=\color{verde}, basicstyle=\ttfamily\scriptsize, numbers=left, numberstyle=\tiny\color{grisos}}
\lstdefinestyle{awl}{basicstyle=\ttfamily\tiny, commentstyle=\color{grisos}\itshape, morecomment=[l]{//},
  keywordstyle=\color{azulg}\bfseries, morekeywords={LD,LDN,A,AN,O,ON,EU,ED,S,R,MOVB,MOVW,MOVD,PLS,TON,CALL,CRET,LPS,LRD,LPP,ALD,OLD,FILL,LDW,LDB,LDD,AW,AD,AB,OW,OD},
  numbers=left, numberstyle=\tiny\color{grisos}, frame=single, backgroundcolor=\color{gris}}
\lstdefinestyle{sql}{language=SQL, keywordstyle=\color{azulg}\bfseries, commentstyle=\color{grisos}\itshape,
  stringstyle=\color{verde}, basicstyle=\ttfamily\tiny, numbers=left, numberstyle=\tiny\color{grisos},
  morekeywords={REFERENCES,SERIAL,SMALLSERIAL,BIGSERIAL,TIMESTAMPTZ,BOOLEAN,INET,RETURNS,TRIGGER,LANGUAGE,DECLARE,RAISE,EXCEPTION,VIEW,INDEX}}

\tikzset{
  ld/.style={line width=0.7pt},
  pics/na/.style={code={\fill[white](-0.12,-0.26) rectangle (0.12,0.26);
     \draw[thick](-0.1,-0.24)--(-0.1,0.24)(0.1,-0.24)--(0.1,0.24);
     \node[font=\scriptsize\ttfamily, above] at (0,0.22){#1};}},
  pics/nc/.style={code={\fill[white](-0.12,-0.26) rectangle (0.12,0.26);
     \draw[thick](-0.1,-0.24)--(-0.1,0.24)(0.1,-0.24)--(0.1,0.24);\draw(-0.17,-0.24)--(0.17,0.24);
     \node[font=\scriptsize\ttfamily, above] at (0,0.22){#1};}},
  pics/pf/.style={code={\fill[white](-0.12,-0.26) rectangle (0.12,0.26);
     \draw[thick](-0.1,-0.24)--(-0.1,0.24)(0.1,-0.24)--(0.1,0.24);
     \node[font=\tiny\bfseries] at (0,0){P};}},
  pics/cmp/.style n args={3}{code={\fill[white](-0.12,-0.26) rectangle (0.12,0.26);
     \draw[thick](-0.1,-0.24)--(-0.1,0.24)(0.1,-0.24)--(0.1,0.24);
     \node[font=\tiny\ttfamily] at (0,0){#2};
     \node[font=\scriptsize\ttfamily, above] at (0,0.22){#1};
     \node[font=\scriptsize\ttfamily, below] at (0,-0.22){#3};}},
  pics/bob/.style n args={3}{code={\fill[white](-0.2,-0.26) rectangle (0.2,0.26);
     \draw[thick](-0.05,0.24) arc (110:250:0.255);\draw[thick](0.05,0.24) arc (70:-70:0.255);
     \node[font=\tiny\bfseries] at (0,0){#2};
     \node[font=\scriptsize\ttfamily, above] at (0,0.22){#1};
     \node[font=\scriptsize\ttfamily, below] at (0,-0.22){#3};}},
  fb/.style={draw, fill=white, font=\scriptsize\ttfamily, align=left, inner sep=3pt, anchor=west},
  barra/.style={line width=2pt},
  estado/.style={draw, rounded corners=3pt, fill=white, align=center, font=\scriptsize, minimum width=22mm, minimum height=10mm, inner sep=3pt},
  estadon/.style={estado, fill=naranja!14, draw=naranja!70!black},
  trans/.style={font=\scriptsize, align=center, fill=white, inner sep=1pt},
  lugar/.style={circle, draw, thick, minimum size=7mm, fill=white},
  transicion/.style={rectangle, fill=black, minimum width=1.6mm, minimum height=7mm, inner sep=0},
  tabla/.style={rectangle split, rectangle split parts=2, draw, fill=white, font=\tiny\ttfamily, align=left, inner sep=2pt, rounded corners=1pt, text width=33mm},
}
\newcommand{\pendiente}[1]{\par\smallskip\noindent\colorbox{naranja!12}{\parbox{\dimexpr\linewidth-2\fboxsep}{\textbf{\color{naranja!80!black}Pendiente.} #1}}\par\smallskip}
\newcommand{\cumple}{\textcolor{verde}{\textbf{cumplido}}}
\newcommand{\parcial}{\textcolor{naranja}{\textbf{parcial}}}
\newcommand{\falta}{\textcolor{rojo}{\textbf{pendiente}}}

\graphicspath{{../../documentacion/planos_isa/pid/}{../../documentacion/planos_isa/electrico/}{../../documentacion/manuales_inversa/}
  {../../database/mockups/}{../../database/documentacion/}{../../database/modelo/}{../../database/interfaz/}
  {../../gemelo-digital/manual/img/}{../../modelo-3d/informe/figuras/}{../../modelo-3d/levantamiento/}{../../../../global/fotos-planta/}{../figuras/}}
\usepackage{pdflscape}
\lstset{showstringspaces=false}
% plano A3 en una página apaisada: \plano[página]{archivo}{leyenda}
\newcommand{\plano}[3][1]{\begin{landscape}\begin{figure}[p]\centering\includegraphics[page=#1, width=\linewidth, height=0.86\textwidth, keepaspectratio]{#2}\caption{#3}\end{figure}\end{landscape}}
\setlength{\headheight}{14pt}
\newcommand{\nlab}{}
\newcommand{\tlab}{}
\lhead{\small Brazo XK-JS1-1 · Informe del laboratorio \nlab}
\rhead{\small UMSA · II/2026}

% \portada{número}{título}{inicio}{entrega}{semanas}
\newcommand{\portada}[5]{%
\renewcommand{\nlab}{#1}%
\begin{titlepage}
\centering
\vspace*{0.6cm}
{\large Universidad Mayor de San Andrés · Facultad de Ingeniería\par}
{\large Carrera de Ingeniería Electrónica\par}
\vspace{0.8cm}
{\color{naranja!80!black}\rule{\linewidth}{1.5pt}\par}
\vspace{0.3cm}
{\Large Informe del Laboratorio Integrado Nº #1\par}
\vspace{0.2cm}
{\huge\bfseries #2\par}
\vspace{0.3cm}
{\color{naranja!80!black}\rule{\linewidth}{1.5pt}\par}
\vspace{0.8cm}
{\large Planta: brazo robótico XK-JS1-1 de 4 grados de libertad\par}
\vspace{0.3cm}
{\large ETN-902 Sistemas de Control II · ETN-1034 Aplicación de Técnicas de Control · ETN-1000 Bases de Datos\par}
\vspace{1cm}
\begin{tabular}{rl}
\textbf{Docente:} & Ing. Jorge Antonio Nava Amador\\
\textbf{Auxiliares:} & Limber Ajnota · Fabián Plata · Franz Choque\\
\textbf{Inicio de trabajo:} & #3\\
\textbf{Entrega:} & #4 (#5)\\
\end{tabular}
\vspace{1cm}

\begin{tabular}{ll}
\multicolumn{2}{l}{\textbf{Escuadra del brazo de 4 GDL}}\\[2pt]
Control (ETN-902 / ETN-1034): & Gabriel Andres Gomez Calla\\
 & Jose Lenin Mamani Choque\\
 & Miguel Alejandro Lima Huarachi\\
 & Román Abner Machicado Arequipa\\[2pt]
Cómputo (ETN-1000): & Miguel Alejandro Lima Huarachi\\
 & Marco Eduardo Chura Rondo\\
\end{tabular}
\vfill
{\small Repositorio: \url{https://github.com/GabAndGom/Brazo-de-4-grados-de-libertad}\par}
{\small La Paz, Bolivia · gestión II/2026\par}
\end{titlepage}
}

% Tabla de entregables (Definition of Done): \begin{dod} ... \end{dod}, filas: entregable & estado & evidencia\\
\newenvironment{dod}{\begin{center}\small\begin{longtable}{>{\raggedright}p{5.2cm} >{\raggedright}p{1.8cm} >{\raggedright\arraybackslash}p{7.6cm}}
\toprule \textbf{Entregable} & \textbf{Estado} & \textbf{Evidencia / observación}\\ \midrule \endhead}
{\bottomrule\end{longtable}\end{center}}

\renewcommand{\tipodoc}{Logros del Laboratorio Integrado}
\renewcommand{\tipocab}{Logros del laboratorio}
\begin{document}
\portada{5}{Gemelo digital, discretización y triggers de seguridad}{martes 29 de septiembre de 2026}{martes 27 de octubre de 2026}{semanas 9 a 13}
\tableofcontents
\newpage
\nota{\textbf{Estado al 18/11/2026.} Laboratorio entregado; cierra la fase 2. Fase 2: modelado dinámico, caja negra y backend (cierre). Este documento explica, objetivo por objetivo, qué pide la guía, qué logró la escuadra, qué es cada cosa y dónde está la evidencia; después explica el trabajo del docente: qué debía revisar en cada hito y qué tiene preparado la escuadra para esa revisión.}
\section{Resumen}
\begin{center}\small\begin{tabularx}{\linewidth}{c >{\raggedright}p{2.4cm} >{\raggedright}X c c}
\toprule \textbf{Nº} & \textbf{Mención} & \textbf{Objetivo} & \textbf{Estado} & \textbf{\%}\\ \midrule
1 & ETN-902 / ETN-1034 & Modelado computacional y simulación (gemelo digital) & \parcial & 60\%\\
2 & ETN-902 / ETN-1034 & Discretización matemática & \parcial & 60\%\\
3 & ETN-1000 & Programación del servidor (lógica de negocios) & \parcial & 90\%\\
4 & ETN-1000 & Seguridad funcional por software & \cumple & 100\%\\
\bottomrule\end{tabularx}\end{center}
El porcentaje mide cuánto del objetivo está hecho y comprobado en la planta o en el repositorio; no cuenta lo que solo está planificado.
\section{Logros respecto de los objetivos}
\subsection{Mención Control Automático (ETN-902 / ETN-1034)}
\subsubsection{Modelado computacional y simulación (gemelo digital)}
\paragraph{Qué pide la guía.} Desarrollar el gemelo digital dinámico de la planta y analizar su respuesta.
\paragraph{Qué se logró.}
\begin{itemize}
\item Gemelo digital propio (Python + navegador) en lugar de Simulink, porque además de simular se conecta al PLC real por PPI.
\item Tres partes: modelo 3D con las piezas de Inventor y la cinemática directa e inversa; PLC simulado que reproduce el programa \codigo{brazo\_gemelo}; y estimador de posición a partir de los pulsos.
\item Validación: estimador frente al PLC con 1,35 mm de error máximo en 170 mm de recorrido (99,2 \%); la cinemática reproduce la pose del ensamble de Inventor a menos de 0,4 mm.
\end{itemize}
\paragraph{Qué es y por qué importa.} Un gemelo digital es una copia en software que se comporta como la planta: permite probar órdenes, límites y secuencias sin arriesgar el hardware. La guía lo usa como filtro: nadie carga código en el PLC sin un gemelo con más de 95 \% de fidelidad.
\paragraph{Evidencia.}
\begin{itemize}\small
\item \codigo{gemelo-digital/app/}
\item \url{https://brazo-4gdl.pages.dev}
\item \codigo{modelo-3d/informe/informe\_cinematica\_modelo\_3d.pdf}
\end{itemize}
\paragraph{Lo que falta.} Medir la posición real de la garra con calibre (índice independiente de la cuenta de pulsos) y, si se exige, un modelo \codigo{.slx} en Simulink.
\paragraph{Estado:} \parcial~(60\%).
\subsubsection{Discretización matemática}
\paragraph{Qué pide la guía.} Pasar el modelo continuo al plano Z y elegir Ts con Nyquist-Shannon y $\tau$/10 $\le$ Ts $\le$ $\tau$/4.
\paragraph{Qué se logró.}
\begin{itemize}
\item Discretización con ZOH: G(z) = b$\cdot$Ts/(z - 1) por eje; en espacio de estados G = I y H = Ts$\cdot$B, es decir x(k+1) = x(k) + Ts$\cdot$B$\cdot$u(k).
\item Ts propuesto por los límites reales del sistema: 25 ms para el control en el PLC (un paquete de 100 pulsos a 4 kHz) y 155 ms para la supervisión desde la PC (una lectura PPI).
\end{itemize}
\paragraph{Qué es y por qué importa.} Discretizar es escribir el modelo como lo ve un controlador digital, que solo actúa cada Ts segundos y mantiene la salida constante entre muestras (retención de orden cero). El Ts debe ser bastante menor que la constante de tiempo $\tau$ de la planta para no perder su dinámica.
\paragraph{Evidencia.}
\begin{itemize}\small
\item \codigo{laboratorios/lab05\_gemelo\_y\_triggers/informe\_lab05.pdf}
\item \codigo{laboratorios/lab05\_gemelo\_y\_triggers/salida\_espacio\_estados.txt}
\end{itemize}
\paragraph{Lo que falta.} Medir $\tau$ con un sensor continuo y comprobar 25 ms $\le$ $\tau$/4; sin $\tau$ la regla de la guía no se puede aplicar.
\paragraph{Estado:} \parcial~(60\%).
\subsection{Mención Sistemas Computacionales (ETN-1000)}
\subsubsection{Programación del servidor (lógica de negocios)}
\paragraph{Qué pide la guía.} Funciones almacenadas y triggers en PL/pgSQL que validen la telemetría.
\paragraph{Qué se logró.}
\begin{itemize}
\item Función \codigo{verificar\_limites\_variables()} y trigger \codigo{muestra\_en\_rango} (BEFORE INSERT OR UPDATE sobre \codigo{muestra}, la tabla de telemetría).
\item Vistas \codigo{v\_ciclo} y la de muestras con tag, y 14 consultas KPI probadas con 18 000 muestras sintéticas.
\end{itemize}
\paragraph{Qué es y por qué importa.} Un trigger es código que la base ejecuta sola cada vez que llega un dato. Así la validación no depende de que el middleware la haga bien: aunque alguien inserte a mano, la regla se cumple.
\paragraph{Evidencia.}
\begin{itemize}\small
\item \codigo{database/scripts/brazo\_4gdl.sql}
\item \codigo{database/scripts/consultas\_kpi.sql}
\end{itemize}
\paragraph{Lo que falta.} Procedimientos almacenados de uso (por ejemplo, abrir y cerrar un ciclo) si el docente los pide.
\paragraph{Estado:} \parcial~(90\%).
\subsubsection{Seguridad funcional por software}
\paragraph{Qué pide la guía.} Detectar excursiones de variables e insertar alertas automáticas en tiempo real.
\paragraph{Qué se logró.}
\begin{itemize}
\item Si un valor sale de su rango (base -263° a 263°, carros 0 a 170 mm, trenes 0 a 20 000 Hz): registra la alarma \codigo{limite\_variable} con \codigo{clock\_timestamp()} (hora UTC con microsegundos), pone \codigo{sistema.en\_operacion = false} y descarta la muestra (RETURN NULL).
\item Pruebas de inyección \codigo{prueba\_reglas.sql}: 7 de 7 correctas (valor fuera de rango, nulo, modo inválido, FK inexistente, tag repetido, borrado con historial).
\end{itemize}
\paragraph{Qué es y por qué importa.} Es la versión en software de un enclavamiento: un dato imposible (por ejemplo, el carro a 250 mm cuando solo recorre 170) indica una falla, así que se aísla, se deja constancia y se marca la planta fuera de operación. Se usa RETURN NULL y no una excepción para que la alerta no se borre con la transacción.
\paragraph{Evidencia.}
\begin{itemize}\small
\item \codigo{database/scripts/prueba\_reglas.sql}
\item \codigo{database/scripts/registro\_ejecucion.txt}
\end{itemize}
\paragraph{Lo que falta.} nada para este objetivo.
\paragraph{Estado:} \cumple~(100\%).
\section{Trabajo del docente}
La guía asigna al docente el papel de auditor: en cada hito revisa el trabajo de la escuadra y decide si cumple antes de pasar a la etapa siguiente. Para cada tarea se explica qué debía revisar, cuándo, y qué tiene la escuadra para presentarle. El porcentaje es el del Excel \codigo{evaluacion\_tareas\_docente.xlsx}.
\subsection{Revisión metodológica de Ts}
\paragraph{Lo que dice la guía.} Evaluar el criterio analítico de selección del periodo de muestreo Ts. \textit{Cuándo:} martes 27 de octubre de 2026.
\paragraph{Qué debía revisar el docente (hitos de la revisión).}
\begin{itemize}
\item[$\square$] Que Ts se justifique con la constante de tiempo $\tau$ de la planta ($\tau$/10 $\le$ Ts $\le$ $\tau$/4) y con Nyquist-Shannon.
\item[$\square$] Que G(z) esté bien calculada con ZOH (40 \% de la defensa junto con la fidelidad).
\end{itemize}
\paragraph{Qué tiene la escuadra para esa revisión.}
\begin{itemize}
\item G(z) = b$\cdot$Ts/(z - 1) y x(k+1) = x(k) + Ts$\cdot$B$\cdot$u(k).
\item Ts de 25 ms (control en el PLC) y 155 ms (supervisión por PPI), justificados por actuación y comunicación, no por $\tau$.
\end{itemize}
\paragraph{Lo que falta para que el docente lo apruebe.} Medir $\tau$ y justificar Ts con $\tau$/10 $\le$ Ts $\le$ $\tau$/4.
\paragraph{Estado:} \parcial~(40\%).
\subsection{Firma del escudo de seguridad}
\paragraph{Lo que dice la guía.} Firmar la conformidad del gemelo digital (más de 95 \% de fidelidad) y habilitar la fase 3. \textit{Cuándo:} martes 27 de octubre de 2026 (cierre de fase 2).
\paragraph{Qué debía revisar el docente (hitos de la revisión).}
\begin{itemize}
\item[$\square$] Índice de ajuste entre el gemelo y la planta real superior al 95 \%, medido con datos independientes.
\item[$\square$] Que el gemelo reproduzca límites, homing y paradas antes de autorizar cargar código en el PLC.
\end{itemize}
\paragraph{Qué tiene la escuadra para esa revisión.}
\begin{itemize}
\item Gemelo 3D con PLC simulado; geometría tomada del modelo de Inventor y estimador con 1,35 mm de error máximo frente al PLC.
\item Demostración en vivo de la app en modo simulado y con la planta.
\end{itemize}
\paragraph{Lo que falta para que el docente lo apruebe.} Medición independiente de la garra, Simulink y firma.
\paragraph{Estado:} \parcial~(50\%).
\subsection{Hito de la fase}
\textbf{Hito de la fase 2 (fin de la semana 13):} gemelos digitales validados en simulación al 95 \% y bases locales creadas, vacías pero probadas con datos sintéticos. La base está cumplida (creada y probada con 18 000 muestras sintéticas). El gemelo existe y funciona, pero su 99,2 \% compara contra la cuenta de pulsos y no contra una medición independiente; falta esa medición y la firma del docente.
\section{Conclusión}
Cumplidos: Seguridad funcional por software. Con avance parcial: Modelado computacional y simulación (gemelo digital), Discretización matemática, Programación del servidor (lógica de negocios).
\par\smallskip Pendientes principales:
\begin{enumerate}
\item Medir la posición real de la garra con calibre (índice independiente de la cuenta de pulsos) y, si se exige, un modelo \codigo{.slx} en Simulink.
\item Medir $\tau$ con un sensor continuo y comprobar 25 ms $\le$ $\tau$/4; sin $\tau$ la regla de la guía no se puede aplicar.
\item Procedimientos almacenados de uso (por ejemplo, abrir y cerrar un ciclo) si el docente los pide.
\end{enumerate}
\par\smallskip\noindent Repositorio: \url{https://github.com/GabAndGom/Brazo-de-4-grados-de-libertad}
\end{document}
